HTTP/1.1 chunked transfer encoding~\cite{rfc9112} is the standard
mechanism by which a client transmits a request body whose length is
unknown when the request line is sent. The specification defines the
wire format but not the granularity at which a parser must deliver
decoded body bytes to the application handler. We measure that
delivery granularity, end to end and under controlled network
conditions, across \textbf{27 production HTTP servers} spanning
\textbf{7 language ecosystems} (Python, Go, Rust, JVM, Node.js,
BEAM, Ruby) and three C servers (\fwid{nginx}, Apache \fwid{httpd},
\fwid{HAProxy}).

We find that \textbf{exactly one server} in our sample
(Microsoft's Kestrel on ASP.NET Core, via the
\fname{System.IO.Pipelines} substrate)
\emph{moves the per-chunk cost off the
\fname{async}/\fname{await} scheduling boundary} on
\textbf{HTTP/1.1 only}: its
\fname{Http1ChunkedEncodingMessageBody.PumpAsync} drains the
socket-readable buffer into a Pipe, loops through every chunk
header present, then wakes the application handler with a single
coalesced \fname{PipeReader.ReadAsync} result. The HTTP/2 path of
the same server (Wave 3) has no pump-equivalent and falls back to
per-DATA-frame dispatch like every other HTTP/2 implementation
we tested. Kestrel still pays a per-chunk parsing cost inside the
H/1 pump (\SI{3.6}{\micro\second}/chunk steady-state), but
eliminates the per-chunk task switch that dominates the other
servers' cost.
Every other server we measure dispatches one parser callback per
chunk and emits one application-receive event per chunk; this
costs between \SI{2.1}{\micro\second} (uvicorn + httptools) and
\SI{114}{\micro\second} (nginx as origin) of single-core CPU per
chunk under our standardized probe.

Four servers exhibit \textbf{superlinear} wall-time scaling
(measured exponents 1.0--2.0 in our $N$ range; mechanistic
arguments support an asymptotic $\Theta(N^2)$): Node's built-in
\fwid{http.Server} (and \fwid{express} and \fwid{fastify} atop it)
via the \fname{process.nextTick} cost of the Readable substrate,
and Elixir's \fwid{bandit} via an \fname{erlang:iolist\_size}/iolist
walk per chunk in \fname{do\_read\_chunk!/4}. A
\SI{1.5}{\mebi\byte} attacker request consumes \emph{measured}
282\,s of Node \fwid{http.Server} CPU at $N{=}\num{250000}$
chunks, and 42\,s of bandit CPU; the cited \fwid{express}
N=\num{250000} cell is extrapolated from a quadratic fit to
smaller cells.

The "deploy behind nginx" mitigation widely advised in framework
documentation is empirically confirmed. With default
\fname{proxy\_request\_buffering~on}, an
upstream-instrumented end-to-end probe (prober
\(\to\) nginx-proxy \(\to\) instrumented Python upstream) shows
the upstream receives \emph{exactly one}
\fname{Content-Length}-framed request with no
\fname{Transfer-Encoding} header and \emph{one} \fname{recv()}
covering both headers and the entire body. With
\fname{proxy\_request\_buffering~off} the same probe forwards
the chunked stream unchanged, with the upstream seeing multiple
\fname{recv()}s and a \fname{Transfer-Encoding: chunked} header.
Apache \fname{mod\_proxy\_http} is documented to behave the
same way (not directly measured). \fwid{HAProxy} does
\emph{not} aggregate: it streams by default, and
\fname{http-buffer-request} only waits for one buffer's worth,
not a per-chunk accumulator. Direct-exposed servers,
HAProxy-fronted deployments, and any deployment with buffering
disabled inherit the full per-chunk cost.

A Wave 3 measurement set extends the survey to HTTP/2 DATA frames
and WebSocket text frames. Every HTTP/2 server in our sample
(6 servers including Kestrel-h2, hypercorn-h2, node-h2, vertx-h2,
go-h2c, rust-hyper-h2) exhibits the same per-frame amplification
shape as HTTP/1 chunked-TE, at 4--10 \(\mu\)s/frame; HTTP/2's
flow control bounds bytes but not frames. WebSocket is the only
protocol in which most ecosystems already batch correctly
(node-ws, gorilla, rust-tungstenite, Kestrel-WS at
0.1--0.7 \(\mu\)s/frame); Python ASGI's WebSocket implementations
(uvicorn + websockets / wsproto) are the at-risk outlier.

We characterize the root cause as a parser/dispatcher boundary
design choice rather than a spec violation, taxonomize the
mitigations available today, and report the responses of upstream
maintainers including the closure of hyper issue
\#4008~\cite{hyper_4008} ("Provide built-in chunked request limits
and CPU-safe streaming helpers") as \fwid{not\_planned} with the
labelled state \fwid{S-waiting-on-author}. We argue that opt-in
aggregation primitives modeled on Kestrel's Pipelines pump should
be exposed by every event-loop HTTP server.
